\documentclass[10pt,a4paper]{amsart}
\usepackage[margin=19mm]{geometry}
\usepackage{amsmath,amssymb,mathtools}
\usepackage[scr=rsfs,cal=euler]{mathalfa}
\usepackage{xfrac}
\usepackage[unicode,hidelinks,bookmarks=false]{hyperref}
\newcommand{\ie}{i.e.\ }
\newcommand{\cf}{cf.\ }
\newcommand{\resp}{respectively\ }
\newcommand{\Subsection}{\subsection}
\providecommand{\nobreakdash}{\nobreak\hbox{-}}
\setlength{\emergencystretch}{2em}
\multlinegap0pt
\usepackage{bm}
\usepackage{bbm}
\usepackage{dsfont}
\usepackage{xspace}
\usepackage{cancel}
\usepackage{mathtools}
\usepackage{amsthm}
\makeatletter
\@ifundefined{theorem}{\newtheorem{theorem}{Theorem}}{}
\makeatother
\def\op{\bigstar}
\def\ra{\rightarrow}
\def\up#1{\mathchoice{\bigl[#1\bigr)}{[#1)}{}{}}
\newcommand{\Cm}[1]{\mathbf{Cm}\, #1}
\newcommand{\Con}[1]{\mathbf{Con}\, #1}
\newcommand{\J}[1]{\mathcal{J}(#1)}
\newcommand{\alg}[1]{\mathbf{#1}}
\newcommand{\cm}[1]{\mathrm{Cm}\, #1}
\newcommand{\deq}{\mathrel{\mathop:}=}
\newcommand{\irr}[1]{\mathrm{I}(#1)}
\newcommand{\var}[1]{\mathsf{#1}}
\title[A generic construction of free algebras in varieties of Hilb]{A generic construction of free algebras in varieties of Hilbert algebras and Brouwerian semilattices}
\author{Tomasz Kowalski, Katarzyna S{\l}omczy\'nska}
\date{Source: arXiv:2607.18852v1 (CC BY 4.0)}
\begin{document}
\maketitle

\noindent\emph{Source and licence.} A generic construction of free algebras in varieties of Hilbert algebras and Brouwerian semilattices. Version arXiv:2607.18852v1 (CC BY 4.0), distributed under the Creative Commons Attribution 4.0 International licence, as recorded in the arXiv licence metadata for this exact version and shown on the abstract page https://arxiv.org/abs/2607.18852v1. Full rights notice: licenses/arxiv-2607.18852-CC-BY-4.0.txt.\par\smallskip

\noindent\textit{Editorial scope.} Counted unit: the Theorem whose source label is \texttt{None} in the article (source0.tex lines 2139--2143, source label \texttt{None}), delivered together with the article's own complete original proof of it (source0.tex lines 2145--2219, one \texttt{proof} environment of 75 lines). No mathematics is written, repaired or re-derived: statement and proof are the article's own text line for line. Required context retained uncounted: thm:PrimeFilter (lines 525--530); thm:irreducible-charact (lines 937--951); eq:Bk (lines 2133--2135). Every definition, display and label of those lines is retained unchanged. Omitted from this package and not redistributed: 1--524, 531--936, 952--2132, 2136--2138, 2144--2144, 2220--2836. Nothing inside a retained excerpt is omitted or altered: every retained body is delivered exactly as the article's lines are written, so no omission receipt of any kind accompanies this package. Citations: the retained excerpts carry no citation command at all, so no bibliography entry of the article is transcribed and no reference is redistributed. Reference adaptation: none was needed and none was made; every reference inside the delivered unit resolves inside it, so no unresolved reference or citation is produced. Printed slips kept literal: no printed word, symbol, hypothesis or number of the counted statement or of its proof was altered. Layout and class: the article's own class is \texttt{amsart} with packages including tikz, eucal, amssymb, mathabx, stmaryrd, algorithm, algpseudocode, enumitem, xcolor; the wrapper is the lane's pinned \texttt{amsart} editorial wrapper at 10pt on A4, which restates the theorem-like environments the retained text needs and adds the article's own macro definitions that the retained text needs (\texttt{\textbackslash Cm}, \texttt{\textbackslash Con}, \texttt{\textbackslash J}, \texttt{\textbackslash alg}, \texttt{\textbackslash cm}, \texttt{\textbackslash deq}, \texttt{\textbackslash irr}, \texttt{\textbackslash op}, \texttt{\textbackslash ra}, \texttt{\textbackslash up}, \texttt{\textbackslash var}), with extra wrapper packages (bm, bbm, dsfont, xspace, cancel, mathtools). The delivered numbering is the unit's own. Assets: no retained span contains a figure environment, a graphics-inclusion command, a PDF-page inclusion, a table or a TikZ picture; no image file is copied into this package, the package declares no asset dependency and no image file is redistributed with it. Licence: version arXiv:2607.18852v1 of this article is distributed under the Creative Commons Attribution 4.0 International licence, as recorded in the arXiv licence metadata for this exact version and shown on the abstract page https://arxiv.org/abs/2607.18852v1. A full rights notice is delivered at licenses/arxiv-2607.18852-CC-BY-4.0.txt. Characters: every retained excerpt is pure ASCII, so the delivered engine, XeLaTeX with its default Latin Modern text font, reports no missing character.
\begin{theorem}\label{thm:PrimeFilter}
Let $\alg{A}$ be a Hilbert algebra or a Brouwerian semilattice.
For any $a\in A$ and any $X\subseteq A$, if $\up{a}\cap X = \emptyset$, then
there is a $\mu\in\Cm{\alg{A}}$ such that $a\in 1/\mu$ and
$X\cap 1/\mu = \emptyset$.
\end{theorem}  

\begin{theorem}\label{thm:irreducible-charact}
Let $\alg{A}\in\var{Hi}$ be finite. Then for any $a\in A\smallsetminus\{1\}$,
the following are equivalent:
\begin{enumerate}
\item $a\in \irr{\alg{A}}$,  
\item $\theta(1,a)\in \J{\Con{\alg{A}}}$,
\item $\forall m\in A: m\ra a\in\{1,a\}$. 
\end{enumerate}
If $\alg{A}\in\var{Br}$ is finite, then the conditions above are
further equivalent to
\begin{enumerate}
\setcounter{enumi}{3}  
\item $a$ is meet-irreducible. 
\end{enumerate}  
\end{theorem}

\begin{equation}\label{eq:Bk}\tag{$B_k$}
\{q_{k,i}(x_1,\dots,x_{k+1})\ra y: 1\leq i\leq k+1\}\ra y = 1.
\end{equation}

\begin{theorem}
Let $\alg{A}\in \var{Hi}$. Then, $\alg{A}\in\var{Hb}_k$ if and only if
the longest antichain in $\Cm{\alg{A}}$ with a common lower bound, is of length
at most $k$.   
\end{theorem}

\begin{proof}
Let $w(\alg{A})$ stand for the length of the longest
antichain in $\Cm{\alg{A}}$ with a common lower bound.
First we show that every 
algebra $\alg{A}$ with $w(\alg{A})\leq k$ satisfies~\eqref{eq:Bk}. 
Suppose $w(\alg{A}) \leq k$, but $\alg{A}$ falsifies~\eqref{eq:Bk}.
Since taking quotients cannot increase the width, we can assume 
$\alg{A}$ is subdirectly irreducible. Take elements
$a_1,\dots,a_{k+1},c\in A$ such that
\begin{equation}\label{eq:fail-q}\tag{\dag}
\alg{A}\models
\{q_{k,i}(a_1,\dots,a_{k+1})\ra c:1\leq i\leq k+1\}\ra c\neq 1.
\end{equation}
We can also assume $c = \op$ (otherwise just take a suitable quotient).
Hence, $q_{k,i}(a_1,\dots,a_{k+1})\ra \op = 1$
for all $i\in\{1,\dots,k+1\}$ and thus $q_{k,i}(a_1,\dots,a_{k+1})\leq\op$
for every $i\in\{1,\dots,k+1\}$.

Next, we claim that the elements $a_i$ and $a_j$ are incomparable, for $i\neq j$.
To see it, suppose that $a_{i_0}\leq a_{j_0}$ for some
distinct $i_0$ and $j_0$. Then
$q_{k,{j_0}}(a_1,\dots,a_{k+1}) = 1$, so 
$q_{k,{j_0}}(a_1,\dots,a_{k+1})\ra\op = \op$, and thus
$$
\alg{A}\models
\{q_{k,i}(a_1,\dots,a_{k+1})\ra \op:1\leq i\leq k+1\}\ra \op = 1,
$$
contradicting~\eqref{eq:fail-q}.
We conclude that $\{a_1,\dots,a_{k+1}\}$ is an antichain.
Therefore, by Theorem~\ref{thm:PrimeFilter},
for each $i\in\{1,\dots,k+1\}$ there exists
$\mu_i\in \Cm{\alg{A}}$ such that $a_i\in 1/\mu_i$ and
$a_j\notin 1/\mu_i$ for all $j\neq i$. Hence,
$\{\mu_i: 1\leq i\leq k+1\}\subseteq\Cm{\alg{A}}$ is an antichain,
and since $\alg{A}$ is subdirectly irreducible, the identity congruence
belongs to $\Cm{\alg{A}}$, so $w(\alg{A}) \deq k+1$ contradicting the supposition.


Conversely, assume $\alg{A}$ satisfies~\eqref{eq:Bk}, yet $w(\alg{A})> k$.  
Assume without loss that $\alg{A}$ is subdirectly irreducible,
and $w(\alg{A}) = k+1$. Let $\mu_1,\dots,\mu_{k+1}\in\cm{\alg{A}}$ be an
antichain with a common lower bound. By CD we get
$\bigcap\{\mu_j: 1\leq j\leq k+1, \ j\neq i\} \not\subseteq \mu_i$,
for every $i\in \{1,\dots,k+1\}$. Therefore,
for each $i$ there is $a_i\in
\bigcap\{\mu_j: 1\leq j\leq k+1, \ j\neq i\}\smallsetminus \mu_i$. We claim that 
$a_1,\dots,a_{k+1}$ is an antichain. To see that,
suppose $a_i\ra a_s = 1$ for some $i\neq s$.
Then $a_i$ and $a_i\ra a_s$ both belong to
$\bigcap\{\mu_j: 1\leq j\leq k+1, \ j\neq i\}$, and so
$a_s\in \bigcap\{\mu_j: 1\leq j\leq k+1, \ j\neq i\}\subseteq \mu_s$.
Hence $a_s\in 1/\mu_s$ contradicting the choice of $a_s$.
Let $\alg{B}$ be the subalgebra of $\alg{A}$  generated by
$\{\op,a_1,\dots,a_{k+1}\}$. Clearly, $\alg{B}$ is finite and subdirectly
irreducible. Now, using Theorem~\ref{thm:PrimeFilter} 
we find $\nu_1,\dots,\nu_{k+1}\in \cm{\alg{B}}$ such that
$\{a_1,\dots,a_{k+1}\}\cap 1/\nu_i = \{a_i\}$. 
Hence $\nu_1,\dots,\nu_{k+1}$ is an antichain. 
We conclude that $\alg{B}$ is finite subdirectly irreducible and
$w(\alg{B}) = k+1$. 

By congruence distributivity and finiteness we have that
$\J{\Con{\alg{B}}}$ and $\Cm{\alg{B}}$ are isomorphic.
Hence, the greatest length of any antichain in $\J{\Con{\alg{B}}}$ is 
$k+1$, and by Theorem~\ref{thm:irreducible-charact} any such antichain of length
$k+1$ in $\J{\Con{\alg{B}}}$ consists of congruences 
$\theta(1,b_1),\dots,\theta(1,b_{k+1})$ such that 
$\{b_i: 1\leq i\leq k+1\}$ is an antichain of irreducible elements. Therefore
$q_{k,i}(b_1,\dots,b_{k+1}) = b_i<\op$, so 
$q_{k,i}(b_1,\dots,b_{k+1})\ra \op = 1$ for each $i\in\{1,\dots,k+1\}$. Hence
$$
\{q_{k,i}(b_1,\dots,b_{k+1})\ra \op:1\leq i\leq k+1\}\ra \op = \op,
$$
contradicting the assumption that $\alg{A}$ satisfies~\eqref{eq:Bk}.
\end{proof}  

\end{document}
